\scanpage{222}
\margin{221}
\begin{center}-- 9 --\end{center}
\margin{n$^\circ$370}
$\mu(f')\ne0$, on a $\mu((D_f-D_{f'})g)=0$ quelle que soit $g\in\mathcal K$ ;
remplaçant $g$ par $(D_f-D_{f'})g$, on obtient
$\mu((D_f-D_{f'})^2g)=0$ quelle que soit $g\in\mathcal K^+$,
donc $D_f=D_{f'}$ puisque le support de $\mu$ est $G$.
Posons donc $D_f=D$. La formule (14) donne
\[
(17)\qquad \mu(f)D(e)=\check\nu(f).
\]
La formule (17) s’étend par linéarité aux $f\in\mathcal K$ telles que $\mu(f)=0$.
On a $D(e)\ne0$ puisque $\check\nu\ne0$.
Ceci établit bien la proportionnalité de $\mu$ et $\check\nu$.

COROLLAIRE. -- \uline{Toute mesure complexe invariante à gauche}
(resp. \uline{à droite}) \uline{sur} $G$
\uline{est proportionnelle à une mesure de Haar à gauche}
(resp. \uline{à droite}).

\uline{Exemple} 1. -- Sur le groupe additif $\mathbb R$, la mesure de Lebesgue $dx$
est une mesure de Haar (Chap. III, § 2, n$^\circ$2, exemple).

\uline{Exemple} 2. -- Pour toute fonction $f\in\mathcal K(\mathbb R_+^*)$,
on a (Fonct. var. réelles, chap. II, § 1, formule (13))
\[
\int_0^{+\infty}\frac{f(x)}x\,dx
=\int_0^{+\infty}\frac{f(tx)}{tx}\,t\,dx
=\int_0^{+\infty}\frac{f(tx)}x\,dx
\]
quel que soit $t>0$ ; la mesure $x^{-1}dx$ est donc une mesure de Haar
sur le groupe multiplicatif $\mathbb R_+^*$.

\uline{Exemple} 3. -- Prenons pour $G$ le tore $\mathbb T=\mathbb R/\mathbb Z$.
Soit $\varphi$ l’application canonique de $\mathbb R$ sur $\mathbb T$.
Pour $f\in\mathcal K(\mathbb T)$, la fonction $f\circ\varphi$ est continue et
périodique de période 1 sur $\mathbb R$, et l’intégrale
$I(f)=\int_a^{a+1}f(\varphi(x))dx$ est indépendante du choix de $a\in\mathbb R$ ;
il est immédiat qu’elle est invariante par translation ; elle définit donc
une mesure de Haar sur $\mathbb T$. Par transport de structure, on en déduit que
$I(f)=\int_a^{a+1}f(e^{2\pi it})dt$ est une mesure de Haar sur le groupe
multiplicatif $\mathbb U$ des nombres complexes de valeur absolue 1
(Top. géné., chap. VIII, § 2, n$^\circ$1).

PROPOSITION 1. -- \uline{Soient} $G$ \uline{un groupe localement compact},
$\mu$ \uline{une mesure de Haar à gauche ou à droite sur} $G$.
\uline{Pour que} $G$ \uline{soit discret}, \uline{il faut et il suffit que}
